\documentclass[11pt]{article}

\usepackage{amsmath, amsthm, amssymb}
\DeclareMathOperator*{\argmax}{arg\,max}
\DeclareMathOperator*{\argmin}{arg\,min}
\usepackage[margin=1in]{geometry}

\theoremstyle{definition}
\newtheorem{definition}{Definition}[section]

\theoremstyle{plain}
\newtheorem{theorem}{Theorem}[section]

\title{}
\date{}

\begin{document}

\section{Utility Maximizing satisfies WARP}
Suppose an agent uses their utility function $u:X \to \mathbb{R}$ to construct a choice correspondence over $X$ by $\forall M \in \mathcal{M}(X)$ setting:
$$C(M)=\argmax_{x \in M} u.$$
Show that $C$ is in fact a choice correspondence and furthermore that $C$ satisfies WARP.

\section{Dutch Book Practice}
Suppose Alice's credences over the events $A, B, A\cap B, A \cup B$ are as follows:
$$P(A)=P(B)=0.25,P(A\cup B)=0.5,P(A\cap B)=0.1.$$
Construct a Dutch book against Alice.

\section{Bayes' Rule Practice} 
Suppose in a population $1\%$ of people have a deadly disease called boobooitis. Thankfully, we have a test that is 99\% accurate in that if someone who has boobooitis takes the test it will be positive 99\% of the time and if someone who doesn't have boobooitis takes the test it will be negative 99\% of the time. Imagine we administer this test to everyone in the population. Of the times the test returns a positive result for someone, how often will that person actually have boobooitis? 

\section{A Consequence of Independence}
Suppose an agent's preferences $\succeq$ over lotteries satisfy independence. Consider two arbitrary lotteries $g, h \in Z(X)$ such that $g \succeq h$. Prove $\forall p,q \in [0,1]$ that if $p>q$ then $g \oplus_p h \succeq g \oplus_q h$.



\section{Valid Representations}
Let $X=\{\text{chocolate}, \text{vanilla}, \text{strawberry}\}$. Consider the utility function $u:X \to \mathbb{R}$ which sets:
$$u(\text{chocolate})=1,u(\text{vanilla})=2,u(\text{strawberry})=3$$
(a) Suppose $u: X \to \mathbb{R}$ is a valid representation of an agent's preferences $\succeq$ over prizes in $X$, i.e.:
$$\forall a, b \in X, a \succeq b \leftrightarrow u(a) \geq u(b).$$
Is $u^3$ also a valid representation of the agent's preferences $\succeq$ over prizes in $X$? Justify your answer using a relevant theorem from lecture. 
\\
\\
(b) Suppose $u: X \to \mathbb{R}$ is a valid representation of an agent's preferences $\succeq$ over lotteries in $Z(X)$, i.e.:
$$\forall g,h \in Z(X), g \succeq h \leftrightarrow \sum_{x \in X}u(x) \cdot g(x) \geq \sum_{x \in X}u(x) \cdot h(x)$$
Is $u^3$ also a valid representation of the agent's preferences $\succeq$ over lotteries in $Z(X)$? Justify your answer using a relevant theorem from lecture. 
\end{document}
